% ============================================================
%  Μοντέλο φθοράς ελαστικού -- μαθηματική διατύπωση
% ============================================================

\label{subsec:wear_model}

Το μοντέλο φθοράς που χρησιμοποιείται στην παρούσα εργασία ακολουθεί την προσέγγιση των
Tremblett και Limebeer \cite{Tremlett2016}, η οποία με τη σειρά της βασίζεται στη θεωρία των δύο
μηχανισμών τριβής ελαστομερούς και επαληθεύεται απο πειραματικά δεδομένα του Grosch \cite{Grosch2004}. Οι τρεις μακροσκοπικοί μηχανισμοί
φθοράς που περιγράφηκαν στο \cref{subsubsec:graining_blistering} -- τριβολογική
φθορά λόγω ολίσθησης, \textit{graining}, και \textit{blistering} -- μοντελοποιούνται ξεχωριστά και υπερτίθενται.
Η μετάβαση μεταξύ της περιοχής χαμηλής και υψηλής θερμοκρασίας μοντελοποιείται με ομαλή συνάρτηση κατ' ανάλογο τρόπο με
τον συντελεστή μείωσης λόγω θερμοκρασίας. 

\subsubsection{Ρυθμός τριβολογικής φθοράς}

Ο ρυθμός τριβολογικής φθοράς $W_p$ συνδέεται άμεσα με το έργο τριβής ολίσθησης στο
πέλμα. Δεδομένου ότι το έργο αυτό μετατρέπεται σχεδόν εξ ολοκλήρου σε θερμότητα (όρος
$Q_1$ του θερμικού μοντέλου), ο ρυθμός φθοράς εκφράζεται μέσω μιας
εκθετικής σχέσης \cite{Grosch2004}:

\begin{equation}
    \dot{w}_p = \dot{w}_{p1} \cdot \left(\frac{Q_1}{Q_{\mathrm{ref}}}\right)^{W_{p2}}
    \label{eq:w_p}
\end{equation}

όπου $Q_{\mathrm{ref}}$ ισχύς αναφοράς στην οποία έγινε η μέτρηση, $\dot{w}_{p1}$ ο ρυθμός φθοράς σε συνθήκες αναφοράς
και $\dot{w}_{p2}$ εκθέτης που ρυθμίζει τη μη γραμμικότητα της απόκρισης. Για $\dot{w}_{p2} > 1$ η
φθορά επιταχύνεται δυσανάλογα σε συνθήκες υψηλής ολίσθησης, γεγονός που επιβεβαιώνεται με
πειραματικές παρατηρήσεις \cite{Grosch2004,Tremlett2016}.

\subsubsection{Ρυθμός graining}

Ο ρυθμός graining $\dot{w}_g$ ενεργοποιείται όταν η θερμοκρασία πέλματος πέφτει κάτω από την
κρίσιμη θερμοκρασία μετάβασης $T_{\mathrm{tp}}$ και μηδενίζεται πάνω από αυτήν.
Μοντελοποιείται με εμπειρική σχέση και παραμέτρους προς διερεύνηση.

\begin{equation}
    \dot{w}_g = W_{g1} \cdot (t_{\mathrm{tp}} - T)^{W_{g2}}
    \label{eq:w_g}
\end{equation}

όπου ο εκθέτης $W_{g2}$ παίζει το ρόλο της επιτάχυνσης της φθοράς μη γραμμικά με τη θερμοκρασία, και ο συντελεστής $W_{g1}$ ελέγχει τη συνολική επίδραση του 
\textit{graining}.

\subsubsection{Ρυθμός blistering}

Αντίστοιχα, ο ρυθμός blistering $\dot{w}_b$ ενεργοποιείται πάνω από την κρίσιμη θερμοκρασία και
μηδενίζεται κάτω από αυτήν. Κατ'αντιστοιχία, ελέγχεται από δύο παραμέτρους, $W_{b1}$ που
καθορίζει την κλίμακα και $W_{b2}$ που ελέγχει την ένταση της εξάρτησης από τη θερμοκρασία.

\begin{equation}
    \dot{w}_b = W_{b1} \cdot (T - t_{\mathrm{tp}})^{W_{b2}}
    \label{eq:Wb}
\end{equation}

Η επιλογή κοινής κρίσιμης θερμοκρασίας $t_{\mathrm{tp}}$ και για τους δύο μηχανισμούς
graining και blistering είναι απλοποίηση: στην πραγματικότητα το θερμικό παράθυρο έχει
μικρό εύρος και όχι μηδενικό, αλλά η επιλογή αυτή διατηρεί τη μαθηματική συμμετρία του
μοντέλου και τεκμηριώνεται στη βιβλιογραφία \cite{West2022,Tremlett2016}.

\subsubsection{Συνολικός ρυθμός φθοράς}

Ο συνολικός ρυθμός φθοράς προκύπτει από την υπέρθεση των τριών μηχανισμών, με ομαλή
μετάβαση στη βέλτιστη θερμοκρασία $t_{tp}$.

\begin{equation}
\dot{w_a}(T_s,Q_1) = (\dot{w}_g+\dot{w}_p) + \Bigl(0.5+0.5sin\bigl(tan^{-1}\bigl(g_s(T_s-t_{tp})\bigr)\bigr)\Bigr)
\cdot \Bigl(\bigl(\dot{w}_b +\dot{w}_p \bigr) -
\bigl(\dot{w}_g +\dot{w}_p \bigr) \Bigr)
    \label{eq:Wtotal}
\end{equation}

Η παράμετρος $g_s$ -- κοινή σε όλες τις ομαλές συναρτήσεις -- ρυθμίζει την καμπυλότητα της μετάβασης:
μικρές τιμές δίνουν ομαλή μετάβαση ενώ όσο αυξάνεται η τιμή της, η τελική μορφή της καμπύλης προσεγγίζει τμηματική συνάρτηση των δύο καμπυλών.

\subsubsection{Ολοκλήρωση σε συσσωρευμένη φθορά}

Ο ρυθμός φθοράς $\dot{w_a}$ έχει διαστάσεις μήκους ανά χρόνο (πάχος πέλματος που χάνεται ανά
δευτερόλεπτο). Η συσσωρευμένη φθορά κατά τη διάρκεια ενός γύρου προκύπτει από
ολοκλήρωση κατά μήκος της τροχιάς, χρησιμοποιώντας τον κανόνα αλυσίδας για μετατροπή από
τη χρονική στη χωρική ανεξάρτητη μεταβλητή:

\begin{equation}
    \frac{dw_a}{ds} = \frac{1}{v}\, \dot{w}_a(T, Q_1)
    \label{eq:dwds}
\end{equation}

με $v$ την ταχύτητα του οχήματος.


Στον παρακάτω πίνακα (\ref{tab:params_wear}) παρατίθενται ενδεικτικές τιμές των παραμέτρων του μοντέλου φθοράς.

\begin{table}[htbp]
    \centering
    \small
    \renewcommand{\arraystretch}{1.2}
    \begin{tabular}{@{}l l l@{}}
        \toprule
        \textbf{Παράμετρος} & \textbf{Τιμή (Μονάδα)} & \textbf{Περιγραφή} \\
        \midrule
        $W_{p1}$             & $0.09$ mm/s                 & Ρυθμός φθοράς ολίσθησης στην ισχύ αναφοράς \\
        $W_{p2}$             & $1.6$                       & Εκθέτης μη γραμμικότητας φθοράς ολίσθησης \\
        $Q_{\mathrm{ref}}$   & $150$ kW                    & Ισχύς αναφοράς όρου φθοράς ολίσθησης\\
        $W_{g1}$             & $4 \cdot 10^{-6}$ mm/K$^2$s & Συντελεστής έντασης graining \\
        $W_{g2}$             & $2$                         & Εκθέτης θερμοκρασιακής εξάρτησης graining \\
        $W_{b1}$             & $8 \cdot 10^{-6}$ mm/K$^2$s & Συντελεστής έντασης blistering \\
        $W_{b2}$             & $2$                         & Εκθέτης θερμοκρασιακής εξάρτησης blistering \\
        $T_{\mathrm{tp}}$    & $100\,^\circ\mathrm{C}$     & Βέλτιστη θερμοκρασία -- μετάβαση graining-blistering \\
        \bottomrule
    \end{tabular}
    \caption{Παράμετροι μοντέλου φθοράς.}
    \label{tab:params_wear}
\end{table}